\documentclass[10pt,a4paper]{amsart}
\usepackage{amsmath,amssymb,amsthm,graphicx}
% Retained minimal preamble for the Stacks Project chapter more-algebra.tex (More on Algebra).
% The chapter file pinned-original-source.tex carries only \input{preamble} on its first line, so no
% preamble text lives inside the chapter itself.  The Stacks Project uses one single shared preamble for
% every chapter; the snapshot stages it identically as sources/stacks-chow/preamble.tex and
% sources/stacks-etale-cohomology/preamble.tex.  The block below is copied mechanically, line for line,
% from sources/stacks-chow/preamble.tex lines 14-18 (the Xy-pic setup) and 154-217 (the native theorem
% declarations together with the mathematical macros).  Nothing was added, removed or redefined; the
% wrapper class and the amsmath/amssymb/amsthm/graphicx base packages are supplied by the assembler.
% Chapter: more-algebra.tex -- More on Algebra.
% For commutative diagrams we use Xy-pic
\usepackage[all]{xy}

% We use 2cell for 2-commutative diagrams.
\xyoption{2cell}

% Theorem environments.
%
\theoremstyle{plain}
\newtheorem{theorem}[subsection]{Theorem}
\newtheorem{proposition}[subsection]{Proposition}
\newtheorem{lemma}[subsection]{Lemma}

\theoremstyle{definition}
\newtheorem{definition}[subsection]{Definition}
\newtheorem{example}[subsection]{Example}
\newtheorem{exercise}[subsection]{Exercise}
\newtheorem{situation}[subsection]{Situation}

\theoremstyle{remark}
\newtheorem{remark}[subsection]{Remark}
\newtheorem{remarks}[subsection]{Remarks}

\numberwithin{equation}{subsection}

% Macros
%
\def\lim{\mathop{\mathrm{lim}}\nolimits}
\def\colim{\mathop{\mathrm{colim}}\nolimits}
\def\Spec{\mathop{\mathrm{Spec}}}
\def\Hom{\mathop{\mathrm{Hom}}\nolimits}
\def\Ext{\mathop{\mathrm{Ext}}\nolimits}
\def\SheafHom{\mathop{\mathcal{H}\!\mathit{om}}\nolimits}
\def\SheafExt{\mathop{\mathcal{E}\!\mathit{xt}}\nolimits}
\def\Sch{\mathit{Sch}}
\def\Mor{\mathop{\mathrm{Mor}}\nolimits}
\def\Ob{\mathop{\mathrm{Ob}}\nolimits}
\def\Sh{\mathop{\mathit{Sh}}\nolimits}
\def\NL{\mathop{N\!L}\nolimits}
\def\CH{\mathop{\mathrm{CH}}\nolimits}
\def\proetale{{pro\text{-}\acute{e}tale}}
\def\etale{{\acute{e}tale}}
\def\QCoh{\mathit{QCoh}}
\def\Ker{\mathop{\mathrm{Ker}}}
\def\Im{\mathop{\mathrm{Im}}}
\def\Coker{\mathop{\mathrm{Coker}}}
\def\Coim{\mathop{\mathrm{Coim}}}

% Boxtimes
%
\DeclareMathSymbol{\boxtimes}{\mathbin}{AMSa}{"02}

%
% Macros for moduli stacks/spaces
%
\def\QCohstack{\mathcal{QC}\!\mathit{oh}}
\def\Cohstack{\mathcal{C}\!\mathit{oh}}
\def\Spacesstack{\mathcal{S}\!\mathit{paces}}
\def\Quotfunctor{\mathrm{Quot}}
\def\Hilbfunctor{\mathrm{Hilb}}
\def\Curvesstack{\mathcal{C}\!\mathit{urves}}
\def\Polarizedstack{\mathcal{P}\!\mathit{olarized}}
\def\Complexesstack{\mathcal{C}\!\mathit{omplexes}}
% \Pic is the operator that assigns to X its picard group, usage \Pic(X)
% \Picardstack_{X/B} denotes the Picard stack of X over B
% \Picardfunctor_{X/B} denotes the Picard functor of X over B
\def\Pic{\mathop{\mathrm{Pic}}\nolimits}
\def\Picardstack{\mathcal{P}\!\mathit{ic}}
\def\Picardfunctor{\mathrm{Pic}}
\def\Deformationcategory{\mathcal{D}\!\mathit{ef}}

% --- appended 2026-10-12 (Chief): the project own definitions of the reference/slogan/history
% display environments, copied mechanically from the shared Stacks preamble lines 8-12
% (\usepackage{verbatim} plus the three \newenvironment lines, which the chapter preamble keeps
% verbatim from the project preamble). Nothing added, removed or redefined: these environments
% carry no mathematical content and the project itself typesets them as comments. Without these
% three lines a counted statement whose own text uses them fails with Environment slogan undefined.
\usepackage{verbatim}
\newenvironment{reference}{\comment}{\endcomment}
\newenvironment{slogan}{\comment}{\endcomment}
\newenvironment{history}{\comment}{\endcomment}
\usepackage{cleveref}
\crefname{originallemma}{original lemma}{original lemmas}
\Crefname{originallemma}{Original lemma}{Original lemmas}
\makeatletter
\newlabel{lemma-lift-complex}{{0BCA}{source}{}{sourcepointer-1}{}}
\newlabel{lemma-lift-complex@cref}{{[originallemma][0][]0BCA}{[1][0][]source}{}{}{}}
\makeatother
\crefname{originallemma}{original lemma}{original lemmas}
\Crefname{originallemma}{Original lemma}{Original lemmas}
\makeatletter
\newlabel{algebra-lemma-NAK}{{00DV}{source}{}{sourcepointer-2}{}}
\newlabel{algebra-lemma-NAK@cref}{{[originallemma][0][]00DV}{[1][0][]source}{}{}{}}
\makeatother
\crefname{originallemma}{original lemma}{original lemmas}
\Crefname{originallemma}{Original lemma}{Original lemmas}
\makeatletter
\newlabel{algebra-lemma-lift-projective-module}{{07LV}{source}{}{sourcepointer-3}{}}
\newlabel{algebra-lemma-lift-projective-module@cref}{{[originallemma][0][]07LV}{[1][0][]source}{}{}{}}
\makeatother
\crefname{originallemma}{original lemma}{original lemmas}
\Crefname{originallemma}{Original lemma}{Original lemmas}
\makeatletter
\newlabel{lemma-check-tor-dimension-modulo-nilpotent}{{0H75}{source}{}{sourcepointer-4}{}}
\newlabel{lemma-check-tor-dimension-modulo-nilpotent@cref}{{[originallemma][0][]0H75}{[1][0][]source}{}{}{}}
\makeatother
\title{Lifting a bounded above complex of projective modules over a nilpotent quotient to the ring itself}
\author{The Stacks Project Authors}
\date{stacks-project:more-algebra.tex:local-snapshot-20261012; GNU FDL 1.2 or later}
\begin{document}
\maketitle
\setcounter{subsection}{1}
\noindent\textit{Editorial note.} All delivered mathematical text is the author's own, copied from the pinned Stacks Project chapter file more-algebra.tex; no statement or proof step was rewritten, reordered or removed. The counted statement is the labeled lemma environment at source lines 19468-19481 (canonical tag 09AR) and the counted proof is the proof environment at source lines 19483-19495, which immediately follows it and is its own proof.\par\medskip
\noindent\textit{Editorial note.} The proof refers to the five properties of the abstract lifting statement by number rather than by content, and each numbered reference is a single sentence whose input is named explicitly; no step of the verification is left out, so the card is not an instance of an omitted argument.\par\medskip
\noindent\textit{Editorial note.} The four references are two to results of this chapter, namely the lifting criterion used as the framework and the boundedness statement, and two to results of other chapters on Nakayama's lemma and on lifting projective modules; the kinds recorded are the type words appearing in the labels themselves.\par\medskip
\noindent\textit{Editorial note.} The preamble span runs to line 90, the current end of the retained preamble; the counted text does not use the appended environments, and the span is taken whole so that the file version is unambiguous.\par\medskip
\noindent\textit{Editorial note.} The source URL is the source-level Stacks Project URL, as required when more than one card is produced from one source; the per-claim tag is carried by claim\_key and by this note, and the counted statement and proof cite no bibliography entry.\par\medskip
\noindent\textit{Allowed standard prerequisites (editorial).}\begin{itemize}
\item The lifting statement for a class of modules closed under the relevant operations applies once five properties hold: closure under projectivity, closure under summands, surjectivity modulo the ideal, lifting of projectives from the quotient and boundedness of the object. invoked at proof lines 19483-19488. Supplies the criterion that is then verified property by property for the class of all projective modules.
\item Nakayama's lemma gives that a map of projective modules which is surjective modulo the ideal is surjective. invoked at proof lines 19486-19487. Establishes one of the five properties for the class of projectives.
\item Every projective module over the quotient lifts to a projective module over the ring, and boundedness of an object can be checked modulo a nilpotent ideal. used at proof lines 19488-19494. Establishes the remaining two properties and hence the conclusion.
\end{itemize}\medskip
\ifdefined\hypertarget\hypertarget{sourcepointer-1}{}\fi
\noindent\textit{Source reference.} Original lemma 0BCA: sources/stacks-more-algebra/pinned-original-source.tex, source lines 19390--19390. The reference is the criterion applied at proof lines 19483-19485; the verification of its five properties, which is the content of this proof, is inside the counted span.\par\medskip
\ifdefined\hypertarget\hypertarget{sourcepointer-2}{}\fi
\noindent\textit{Source reference.} Original lemma 00DV: sources/stacks-chow/tags.txt, source lines 512--512. Nakayama's lemma is invoked at proof lines 19486-19487 to verify that a map of projectives surjective modulo the ideal is surjective; it is one of the five property checks inside the counted span.\par\medskip
\ifdefined\hypertarget\hypertarget{sourcepointer-3}{}\fi
\noindent\textit{Source reference.} Original lemma 07LV: sources/stacks-chow/tags.txt, source lines 9705--9705. The reference supplies the lifting of projective modules over the quotient, used at proof lines 19488-19490 to verify the fourth property; the use made of it is inside the counted span.\par\medskip
\ifdefined\hypertarget\hypertarget{sourcepointer-4}{}\fi
\noindent\textit{Source reference.} Original lemma 0H75: sources/stacks-more-algebra/pinned-original-source.tex, source lines 17307--17307. The reference is applied at proof lines 19491-19494 to see that the object is bounded below, which is the last of the five properties; the conclusion drawn from it is inside the counted span.\par\medskip
\begin{lemma}
\label{lemma-lift-complex-projectives}
Let $R$ be a ring. Let $I \subset R$ be an ideal. Let $E^\bullet$
be a complex of $R/I$-modules. Let $K$ be an object of $D(R)$. Assume that
\begin{enumerate}
\item $E^\bullet$ is a bounded above complex of projective $R/I$-modules,
\item $K \otimes_R^\mathbf{L} R/I$ is represented by $E^\bullet$ in
$D(R/I)$, and
\item $I$ is a nilpotent ideal.
\end{enumerate}
Then there exists a bounded above complex $P^\bullet$ of projective
$R$-modules representing $K$ in $D(R)$ such that $P^\bullet \otimes_R R/I$
is isomorphic to $E^\bullet$.
\end{lemma}

\begin{proof}
We apply Lemma \ref{lemma-lift-complex} using the class $\mathcal{P}$
of all projective $R$-modules. Properties (1) and (2) of the lemma
are immediate. Property (3) follows from Nakayama's lemma
(Algebra, Lemma \ref{algebra-lemma-NAK}).
Property (4) follows from the fact that we can lift projective
$R/I$-modules to projective $R$-modules, see
Algebra, Lemma \ref{algebra-lemma-lift-projective-module}.
To see that (5) holds it suffices to show that $K$ is in $D^{-}(R)$.
Since we are given that $K \otimes_R^\mathbf{L} R/I$ is in $D^{-}(R/I)$
because $E^\bullet$ is bounded above, this follows from
Lemma \ref{lemma-check-tor-dimension-modulo-nilpotent}.
\end{proof}

\end{document}
